\documentclass[12pt]{amsart}

\usepackage{amsmath}
\usepackage{amssymb}
\usepackage{amsthm}
\usepackage{mathtools}
\usepackage[colorlinks=true,citecolor=blue,urlcolor=blue,linkcolor=blue]{hyperref}

% Theorem environments
\newtheorem{theorem}{Theorem}[section]
\newtheorem{lemma}[theorem]{Lemma}
\newtheorem{proposition}[theorem]{Proposition}
\newtheorem{corollary}[theorem]{Corollary}
\newtheorem{conjecture}[theorem]{Conjecture}

\theoremstyle{definition}
\newtheorem{definition}[theorem]{Definition}
\newtheorem{assumption}[theorem]{Assumption}
\newtheorem{example}[theorem]{Example}

\theoremstyle{remark}
\newtheorem{remark}[theorem]{Remark}

% Operators
\DeclareMathOperator{\tr}{tr}
\DeclareMathOperator{\diag}{diag}
\DeclareMathOperator{\rank}{rank}
\DeclareMathOperator{\MSE}{MSE}
\DeclareMathOperator{\relu}{ReLU}

% Shortcuts
\newcommand{\R}{\mathbb{R}}
\newcommand{\E}{\mathbb{E}}
\newcommand{\Prob}{\mathbb{P}}
\newcommand{\bx}{\mathbf{x}}
\newcommand{\bX}{\mathbf{X}}
\newcommand{\bJ}{\mathbf{J}}
\newcommand{\bM}{\mathbf{M}}
\newcommand{\bU}{\mathbf{U}}
\newcommand{\bV}{\mathbf{V}}
\newcommand{\bI}{\mathbf{I}}
\newcommand{\bu}{\mathbf{u}}
\newcommand{\bv}{\mathbf{v}}
\newcommand{\bw}{\mathbf{w}}
\newcommand{\ba}{\mathbf{a}}
\newcommand{\bg}{\mathbf{g}}
\newcommand{\norm}[1]{\left\lVert #1 \right\rVert}
\newcommand{\abs}[1]{\left\lvert #1 \right\rvert}
\newcommand{\ip}[2]{\langle #1, #2 \rangle}

\title[Noise Propagation and Benign Overfitting]{%
  Noise Propagation Operators and a Two-Parameter\\
  Characterization of Benign Overfitting}

\author{C$\Lambda$ / Lightman}
\address{}
\email{Lightman.chang@gmail.com}

\date{\today}

\subjclass[2020]{62J05, 68T07, 62G20}

\keywords{Benign overfitting, overparameterized models, generalization bounds,
  noise propagation, Jacobian spectrum, neural tangent kernel}

\begin{document}

\begin{abstract}
Overparameterized models that interpolate noisy training data can nonetheless
generalize well, a phenomenon known as benign overfitting.  Existing taxonomies,
notably that of Mallinar et al.\ (2022), characterize the
benign--tempered--catastrophic trichotomy through a single eigenvalue sequence of
the kernel matrix.  We introduce the \emph{noise propagation operator}
$\bM = \boldsymbol{\Sigma}^{-1}\boldsymbol{\Gamma}\boldsymbol{\Sigma}^{-1}$,
where $\boldsymbol{\Sigma}$ collects the singular values of the training Jacobian
and $\boldsymbol{\Gamma}$ measures the test--train cross-correlation of gradient
features.  This operator yields a unified generalization bound whose excess risk
decomposes into a signal-bias term, a noise-variance term $\sigma^{2}\tr(\bM)$,
and a concentration tail controlled by the Frobenius norm of $\bM$.  Under power-law
decay assumptions $\sigma_{j}\sim j^{-\alpha}$ and
$\rho_{j}\sim j^{-\beta}$, we prove a sharp dichotomy: the noise-variance term
$\sigma^{2}\tr(\bM)$ is finite if and only if
$\beta > \alpha + \tfrac{1}{2}$, diverges logarithmically when
$\beta = \alpha + \tfrac{1}{2}$, and diverges polynomially otherwise.  The two-parameter
$(\alpha,\beta)$ framework separates the mechanism by which noise enters the model
(governed by $\alpha$) from the mechanism by which noise is expressed at test
points (governed by $\beta$), thereby explaining why feature learning can convert
catastrophic overfitting into benign overfitting by increasing $\beta$ while
$\alpha$ remains fixed.  We illustrate the framework with a spectral-separation
result for two-layer ReLU teacher--student networks.
\end{abstract}

\maketitle

%% ============================================================
\section{Introduction}\label{sec:intro}
%% ============================================================

A central puzzle in the theory of overparameterized learning is that models with
far more parameters than training samples can interpolate noisy data---fitting
every label exactly---yet still predict accurately on unseen inputs.  Classical
bias--variance trade-off reasoning predicts that such interpolation should
amplify noise and destroy generalization, but modern neural networks routinely
defy this prediction.  Understanding when and why interpolation is compatible
with generalization is one of the most active questions in statistical learning
theory.

\medskip
\noindent\textbf{What problem do we address?}
We seek a quantitative criterion that determines whether an interpolating
predictor generalizes well (benign overfitting), moderately (tempered
overfitting), or poorly (catastrophic overfitting), in terms of the spectral
geometry of the model's Jacobian and its interaction with the test distribution.

\medskip
\noindent\textbf{Why is this important?}
Existing criteria for benign overfitting, while powerful, rely on a single
spectral sequence.  In fixed-kernel regression the kernel eigenvalues
simultaneously control how noise is absorbed during training and how it is
expressed at test time, so a single sequence suffices.  However, after feature
learning the training-time and test-time spectral behaviors can decouple: the
learned features may suppress noise in directions that are invisible at test
points even if those directions carry substantial noise energy in the training
fit.  A single-parameter criterion cannot capture this decoupling.

\medskip
\noindent\textbf{What have others done?}
Bartlett, Long, Lugosi, and Tsigler~\cite{BLLT2020} established the first
rigorous benign-overfitting result for linear regression, showing that the
minimum-norm interpolant generalizes when the effective rank of the covariance
tail is large relative to the sample size.  Tsigler and
Bartlett~\cite{TB2023} extended these results to ridge regression with sharp
non-asymptotic bounds.  Mallinar et al.~\cite{MSCBR2022} introduced a
taxonomy---benign, tempered, and catastrophic---based on the eigenvalue decay
rate of the kernel matrix, providing a unified qualitative picture across
kernel methods and neural networks.  On the optimization side, Jacot, Gabriel,
and Hongler~\cite{JGH2018} introduced the Neural Tangent Kernel (NTK), which
linearizes the network around initialization and connects overparameterized
training dynamics to kernel regression.  Hastie, Montanari, Rosset, and
Tibshirani~\cite{HMRT2022} gave precise asymptotic characterizations of the
double-descent phenomenon in least-squares regression.  Boursier and
Flammarion~\cite{BF2022} analyzed gradient flow in two-layer ReLU
teacher--student networks, proving that the learned weights align with the
teacher directions.  Belkin, Hsu, Ma, and Mandal~\cite{BHMM2019} empirically
demonstrated the double-descent curve.  Mei and Montanari~\cite{MM2022}
analyzed generalization in random feature models, providing sharp asymptotics
that connect kernel regression to neural network behaviour.  Rakhlin and
Zhai~\cite{RZ2019} gave minimax-optimal risk bounds for kernel regression
that highlight the role of the effective dimension.

\medskip
\noindent\textbf{What is new here?}
We introduce the \emph{noise propagation operator}
$\bM = \boldsymbol{\Sigma}^{-1}\boldsymbol{\Gamma}\boldsymbol{\Sigma}^{-1}$,
which factors the test-time noise amplification into two independent
components:
\begin{itemize}
  \item $\boldsymbol{\Sigma}$: the singular values of the training Jacobian,
    governing how noise enters the interpolating solution (decay rate $\alpha$);
  \item $\boldsymbol{\Gamma}$: the test--train cross-correlation of gradient
    features, governing how noise is expressed at test points (decay rate
    $\beta$).
\end{itemize}
In fixed-kernel regression, $\boldsymbol{\Gamma}$ is determined by
$\boldsymbol{\Sigma}$ (both are controlled by the same kernel eigenvalues), so
the framework reduces to a single-parameter condition.  After feature learning,
$\boldsymbol{\Gamma}$ can change independently of $\boldsymbol{\Sigma}$: the
learned representation can increase $\beta$ (making noise directions invisible
at test time) while $\alpha$ remains fixed.  This separation provides a
mechanistic explanation for why feature learning improves generalization beyond
what any fixed-kernel theory can predict.

Our main contributions are:
\begin{enumerate}
  \item A \textbf{unified generalization bound} (Theorem~\ref{thm:main})
    expressing the test mean squared error as the sum of a signal-bias term, the
    noise-variance term $\sigma^{2}\tr(\bM)$, and a concentration tail
    controlled by $\norm{\bM}_{F}$.
  \item A \textbf{benign-overfitting dichotomy} (Theorem~\ref{thm:dichotomy})
    giving a sharp two-parameter condition: overfitting is benign if and only if
    $\beta > \alpha + \tfrac{1}{2}$.
  \item An \textbf{application to two-layer ReLU networks}
    (Proposition~\ref{prop:teacher-student}) demonstrating that gradient-flow
    training achieves spectral separation in the Jacobian, yielding benign
    generalization.
\end{enumerate}

%% ============================================================
\section{Preliminaries}\label{sec:prelim}
%% ============================================================

\subsection{Setting and notation}

Let $S = \{(\bx_{i}, y_{i})\}_{i=1}^{n}$ be a training set with
$\bx_{i} \in \R^{d}$ and $y_{i} \in \R$.  We write
$\mathbf{y} = (y_{1}, \ldots, y_{n})^{\top} \in \R^{n}$ for the label vector.
We assume a standard regression model
\begin{equation}\label{eq:data-model}
  y_{i} = g^{*}(\bx_{i}) + \xi_{i}, \qquad i = 1, \ldots, n,
\end{equation}
where $g^{*}\colon \R^{d} \to \R$ is the unknown ground-truth function and
$\xi_{1}, \ldots, \xi_{n}$ are independent, mean-zero, sub-Gaussian random
variables with variance proxy $\sigma^{2}$.  That is, for every
$\lambda \in \R$,
\begin{equation}\label{eq:sub-gaussian}
  \E\bigl[e^{\lambda \xi_{i}}\bigr] \leq e^{\lambda^{2}\sigma^{2}/2}.
\end{equation}

\begin{remark}[Notation convention]\label{rem:notation}
Throughout this paper, $\sigma$ (without subscript) denotes the noise standard
deviation, while $\sigma_{j}$ (with subscript) denotes the $j$-th singular value
of the Jacobian $\bJ$.  The diagonal matrix of singular values is
$\boldsymbol{\Sigma} = \diag(\sigma_{1}, \ldots, \sigma_{n})$.
\end{remark}

Let $f(\cdot\,; \theta)\colon \R^{d} \to \R$ be a parametric predictor with
parameter vector $\theta \in \R^{p}$.  We assume that $\theta^{*}$ is the \emph{minimum-norm} interpolating
parameter vector relative to a reference $\theta_{0}$:
\begin{equation}\label{eq:interpolation}
  f(\bx_{i}; \theta^{*}) = y_{i}, \qquad i = 1, \ldots, n,
\end{equation}
with $\theta^{*} - \theta_{0} = \bV\boldsymbol{\Sigma}^{-1}\bU^{\top}
(\mathbf{y} - f(\bX; \theta_{0}))$ (the minimum-norm linearized solution).
For simplicity we take $\theta_{0}$ such that $f(\bx_{i}; \theta_{0}) = 0$
for all $i$, which can always be achieved by redefining $f$.

\subsection{Jacobian and its singular value decomposition}

\begin{definition}[Training Jacobian]\label{def:jacobian}
The \emph{training Jacobian} at $\theta^{*}$ is the matrix
$\bJ \in \R^{n \times p}$ whose $i$-th row is the gradient of the prediction at
the $i$-th training point:
\begin{equation}\label{eq:jacobian}
  \bJ_{i,:} = \nabla_{\theta} f(\bx_{i}; \theta^{*})^{\top}, \qquad
  i = 1, \ldots, n.
\end{equation}
\end{definition}

Let the singular value decomposition (SVD) of $\bJ$ be
\begin{equation}\label{eq:svd}
  \bJ = \bU \boldsymbol{\Sigma} \bV^{\top},
\end{equation}
where $\bU = [\bu_{1}, \ldots, \bu_{n}] \in \R^{n \times n}$ is orthogonal,
$\bV = [\bv_{1}, \ldots, \bv_{p}] \in \R^{p \times p}$ is orthogonal (with
only the first $n$ columns relevant when $p > n$), and
$\boldsymbol{\Sigma} = \diag(\sigma_{1}, \ldots, \sigma_{n})$ with
$\sigma_{1} \geq \sigma_{2} \geq \cdots \geq \sigma_{n} > 0$.  The assumption
$\sigma_{n} > 0$ is equivalent to requiring that $\bJ$ has full row rank, which
is a necessary condition for interpolation via linearized perturbation of the
parameters.

\subsection{Gradient feature functions and cross-correlation}

\begin{definition}[Gradient feature functions]\label{def:psi}
For each $j = 1, \ldots, n$, the \emph{$j$-th gradient feature function} is
\begin{equation}\label{eq:psi}
  \psi_{j}(\bx) = \nabla_{\theta} f(\bx; \theta^{*})^{\top} \bv_{j},
\end{equation}
where $\bv_{j}$ is the $j$-th right singular vector of $\bJ$.  Evaluated at a
test point $\bx$, the scalar $\psi_{j}(\bx)$ measures the projection of the
test gradient onto the $j$-th training principal direction.
\end{definition}

\begin{definition}[Cross-correlation matrix]\label{def:Gamma}
The \emph{cross-correlation matrix} $\boldsymbol{\Gamma} \in \R^{n \times n}$ is
defined by
\begin{equation}\label{eq:Gamma}
  \boldsymbol{\Gamma}_{jl}
  = \E_{\bx}\bigl[\psi_{j}(\bx)\,\psi_{l}(\bx)\bigr],
  \qquad j, l = 1, \ldots, n,
\end{equation}
where the expectation is over the test distribution.  The matrix
$\boldsymbol{\Gamma}$ is symmetric and positive semidefinite.  We write
$\rho_{j} = \sqrt{\boldsymbol{\Gamma}_{jj}}$ for the root-mean-square magnitude
of the $j$-th gradient feature at test points.
\end{definition}

\subsection{The noise propagation operator}

\begin{definition}[Noise propagation operator]\label{def:M}
The \emph{noise propagation operator} is the matrix
\begin{equation}\label{eq:M}
  \bM = \boldsymbol{\Sigma}^{-1} \boldsymbol{\Gamma}\,
        \boldsymbol{\Sigma}^{-1} \in \R^{n \times n}.
\end{equation}
Since $\boldsymbol{\Gamma}$ is positive semidefinite and
$\boldsymbol{\Sigma}^{-1}$ is positive diagonal, $\bM$ is symmetric and positive
semidefinite.
\end{definition}

\begin{remark}\label{rem:M-interpretation}
The operator $\bM$ captures the full pipeline of noise amplification: the noise
vector $\boldsymbol{\xi} \in \R^{n}$ is first projected onto the training
singular directions (the columns of $\bU$), scaled by $\sigma_{j}^{-1}$ (the
inverse singular values of $\bJ$), and then the resulting coefficients are
propagated to test points through the cross-correlation $\boldsymbol{\Gamma}$.
The factor $\boldsymbol{\Sigma}^{-1}$ appears twice because the interpolation
condition forces the noise coefficient in direction $j$ to be proportional to
$\sigma_{j}^{-1}$, and the test-time variance of that coefficient depends on its
squared magnitude $\sigma_{j}^{-2}$.
\end{remark}

\subsection{Linearization residual}

Our analysis relies on a local linearization of the predictor around a reference
parameter $\theta_{0}$ (for example, the initialization).

\begin{assumption}[Linearization]\label{ass:lin}
There exists a constant $\delta_{\mathrm{lin}} \geq 0$ such that for every test
point $\bx$ in the support of the test distribution,
\begin{equation}\label{eq:lin-residual}
  \bigl\lvert
    f(\bx; \theta^{*}) - f(\bx; \theta_{0})
    - \nabla_{\theta} f(\bx; \theta_{0})^{\top}(\theta^{*} - \theta_{0})
  \bigr\rvert
  \leq \delta_{\mathrm{lin}}.
\end{equation}
\end{assumption}

When $\delta_{\mathrm{lin}} = 0$ the predictor is exactly linear in the
parameters (as in kernel regression or the NTK regime).  For finite-width
networks, $\delta_{\mathrm{lin}}$ quantifies the quality of the first-order
Taylor approximation.

\subsection{Signal bias}

\begin{definition}[Signal bias]\label{def:signal-bias}
The \emph{signal bias} is defined as
\begin{equation}\label{eq:signal-bias}
  B_{\mathrm{signal}}^{2}
  = \E_{\bx}\Biggl[
      \Biggl(
        \sum_{j=1}^{n}
        \frac{\bu_{j}^{\top} \bg^{*}}{\sigma_{j}}\,\psi_{j}(\bx)
        - g^{*}(\bx)
      \Biggr)^{\!2}
    \Biggr],
\end{equation}
where $\bg^{*} = (g^{*}(\bx_{1}), \ldots, g^{*}(\bx_{n}))^{\top} \in \R^{n}$
is the vector of noiseless labels.  This term measures the approximation error
incurred by the linearized interpolant in reconstructing the true function, and
depends on the richness of the feature space but not on the noise.
\end{definition}

%% ============================================================
\section{Main Results}\label{sec:main}
%% ============================================================

\subsection{Unified generalization bound}

\begin{theorem}[Unified Generalization Bound]\label{thm:main}
Let $S = \{(\bx_{i}, y_{i})\}_{i=1}^{n}$ be generated according to the
model~\eqref{eq:data-model} with sub-Gaussian noise
satisfying~\eqref{eq:sub-gaussian}.  Let $\theta^{*}$ be an interpolating
solution satisfying~\eqref{eq:interpolation}, and let $\bJ$,
$\boldsymbol{\Sigma}$, $\bV$, $\psi_{j}$, $\boldsymbol{\Gamma}$, and $\bM$ be
as in Definitions~\ref{def:jacobian}--\ref{def:M}.  Suppose
Assumption~\ref{ass:lin} holds with residual $\delta_{\mathrm{lin}}$.  Then
there exists an absolute constant $C > 0$ (depending only on the
sub-Gaussian norm of $\xi_{i}$) such that for every
$\delta \in (0, 1)$, with probability at least $1 - \delta$ over the noise
$\boldsymbol{\xi} = (\xi_{1}, \ldots, \xi_{n})^{\top}$,
\begin{equation}\label{eq:main-bound}
  \MSE_{\mathrm{test}}
  \leq C\bigl(B_{\mathrm{signal}}^{2}
       + \sigma^{2}\,\tr(\bM)
       + \sigma^{2}\,\norm{\bM}_{F}\,\sqrt{\log(1/\delta)}
       + \delta_{\mathrm{lin}}^{2}\bigr),
\end{equation}
where
$\MSE_{\mathrm{test}} = \E_{\bx}\bigl[(f(\bx; \theta^{*}) - g^{*}(\bx))^{2}\bigr]$
is the test mean squared error (conditional on the training data and noise
realization), and $\norm{\bM}_{F} = (\sum_{j,l} \bM_{jl}^{2})^{1/2}$ is the
Frobenius norm.
\end{theorem}

\begin{proof}
The proof proceeds in four steps.

\medskip
\noindent\textbf{Step 1: Decomposition of the test prediction error.}
By the interpolation condition~\eqref{eq:interpolation} and the data
model~\eqref{eq:data-model}, the parameter perturbation
$\theta^{*} - \theta_{0}$ is determined by the system
$\bJ(\theta^{*} - \theta_{0}) \approx \mathbf{y} - f(\bX; \theta_{0})$, where
the approximation is controlled by the linearization residual.  Under
Assumption~\ref{ass:lin}, the test prediction error admits the decomposition
\begin{equation}\label{eq:decomp}
  f(\bx; \theta^{*}) - g^{*}(\bx) = S(\bx) + N(\bx) + R(\bx),
\end{equation}
where the three terms are defined as follows.

The \emph{signal reconstruction term} is
\begin{equation}\label{eq:signal-term}
  S(\bx) = \sum_{j=1}^{n} \frac{\bu_{j}^{\top} \bg^{*}}{\sigma_{j}}\,
            \psi_{j}(\bx) - g^{*}(\bx).
\end{equation}
This term captures how well the linearized interpolant reconstructs the noiseless
function $g^{*}$.

The \emph{noise leakage term} is
\begin{equation}\label{eq:noise-term}
  N(\bx) = \sum_{j=1}^{n} \frac{\bu_{j}^{\top} \boldsymbol{\xi}}{\sigma_{j}}\,
            \psi_{j}(\bx).
\end{equation}
This term captures how training noise propagates to the test prediction.

The \emph{linearization residual term} satisfies $\abs{R(\bx)} \leq \delta_{\mathrm{lin}}$ by Assumption~\ref{ass:lin}.

To verify this decomposition, note that the interpolation condition gives
$\bJ(\theta^{*} - \theta_{0}) = \bg^{*} + \boldsymbol{\xi} - f(\bX; \theta_{0}) + \mathbf{r}$
where $\mathbf{r}$ collects training-point residuals.  Using the SVD
$\bJ = \bU\boldsymbol{\Sigma}\bV^{\top}$, the minimum-norm solution is
$\theta^{*} - \theta_{0} = \bV\boldsymbol{\Sigma}^{-1}\bU^{\top}(\bg^{*} + \boldsymbol{\xi} - f(\bX;\theta_{0}) + \mathbf{r})$.
The linearized test prediction is
\begin{align*}
  f(\bx; \theta^{*})
  &\approx f(\bx; \theta_{0})
    + \nabla_{\theta} f(\bx; \theta_{0})^{\top}(\theta^{*} - \theta_{0}) \\
  &= f(\bx; \theta_{0})
    + \sum_{j=1}^{n} \frac{\bu_{j}^{\top}(\bg^{*} + \boldsymbol{\xi}
      - f(\bX;\theta_{0}) + \mathbf{r})}{\sigma_{j}}\,\psi_{j}(\bx).
\end{align*}
Subtracting $g^{*}(\bx)$ and rearranging yields~\eqref{eq:decomp}, where the
signal reconstruction involves $\bu_{j}^{\top}\bg^{*}/\sigma_{j}$, the noise
leakage involves $\bu_{j}^{\top}\boldsymbol{\xi}/\sigma_{j}$, and all
remaining terms (including those involving $f(\bX;\theta_{0})$ and
$\mathbf{r}$) are absorbed into $S(\bx)$ and $R(\bx)$.

\medskip
\noindent\textbf{Step 2: Expected noise variance equals $\sigma^{2}\tr(\bM)$.}
We compute the expected test MSE contribution from the noise term.  Defining the
random variable
\begin{equation}\label{eq:Z}
  Z = \E_{\bx}\bigl[N(\bx)^{2}\bigr]
    = \E_{\bx}\Biggl[
        \Biggl(
          \sum_{j=1}^{n}
          \frac{\bu_{j}^{\top}\boldsymbol{\xi}}{\sigma_{j}}\,\psi_{j}(\bx)
        \Biggr)^{\!2}
      \Biggr],
\end{equation}
we expand the square and exchange the expectation over $\bx$ with the sum:
\begin{align}
  Z &= \sum_{j=1}^{n} \sum_{l=1}^{n}
       \frac{\bu_{j}^{\top}\boldsymbol{\xi}}{\sigma_{j}}\,
       \frac{\bu_{l}^{\top}\boldsymbol{\xi}}{\sigma_{l}}\,
       \E_{\bx}\bigl[\psi_{j}(\bx)\,\psi_{l}(\bx)\bigr] \notag \\
    &= \sum_{j=1}^{n} \sum_{l=1}^{n}
       \frac{(\bu_{j}^{\top}\boldsymbol{\xi})\,
             (\bu_{l}^{\top}\boldsymbol{\xi})}
            {\sigma_{j}\,\sigma_{l}}\,
       \boldsymbol{\Gamma}_{jl} \notag \\
    &= \boldsymbol{\xi}^{\top}
       \bU \boldsymbol{\Sigma}^{-1} \boldsymbol{\Gamma}\,
       \boldsymbol{\Sigma}^{-1} \bU^{\top}
       \boldsymbol{\xi} \notag \\
    &= \boldsymbol{\xi}^{\top}
       \bU \bM \bU^{\top}
       \boldsymbol{\xi}. \label{eq:Z-quadratic}
\end{align}
Taking the expectation over $\boldsymbol{\xi}$, using
$\E[\boldsymbol{\xi}\boldsymbol{\xi}^{\top}] = \sigma^{2}\bI_{n}$ and the
orthogonality of $\bU$:
\begin{align}
  \E_{\boldsymbol{\xi}}[Z]
  &= \sigma^{2}\,\tr\bigl(\bU \bM \bU^{\top}\bigr) \notag \\
  &= \sigma^{2}\,\tr\bigl(\bU^{\top} \bU \bM\bigr) \notag \\
  &= \sigma^{2}\,\tr(\bM). \label{eq:EZ}
\end{align}
This establishes that the expected noise contribution to the test MSE is
$\sigma^{2}\tr(\bM)$.

\medskip
\noindent\textbf{Step 3: Concentration via the Hanson--Wright inequality.}
The quadratic form $Z = \boldsymbol{\xi}^{\top}\bU\bM\bU^{\top}\boldsymbol{\xi}$
from~\eqref{eq:Z-quadratic} involves a sub-Gaussian random vector
$\boldsymbol{\xi}$ and the positive semidefinite matrix
$\mathbf{A} = \bU\bM\bU^{\top}$.  Note that
$\norm{\mathbf{A}}_{F} = \norm{\bM}_{F}$ and
$\norm{\mathbf{A}}_{\mathrm{op}} = \norm{\bM}_{\mathrm{op}}$ since $\bU$ is
orthogonal.

By the Hanson--Wright inequality (see, e.g., Rudelson and
Vershynin~\cite{RV2013}), there exists an absolute constant $c > 0$ such that
for all $t > 0$,
\begin{equation}\label{eq:HW}
  \Prob\bigl[\abs{Z - \E[Z]} > t\bigr]
  \leq 2\exp\biggl(
    -c\,\min\biggl(
      \frac{t^{2}}{\sigma^{4}\norm{\bM}_{F}^{2}},\;
      \frac{t}{\sigma^{2}\norm{\bM}_{\mathrm{op}}}
    \biggr)
  \biggr).
\end{equation}
Setting the right-hand side equal to $\delta$ and solving for $t$ in the
Frobenius-norm regime (which dominates when $t$ is not too large), we obtain that
with probability at least $1 - \delta$,
\begin{equation}\label{eq:Z-tail}
  Z \leq \E[Z]
    + C_{1}\,\sigma^{2}\,\norm{\bM}_{F}\,\sqrt{\log(1/\delta)},
\end{equation}
where $C_{1} > 0$ is an absolute constant depending only on $c$ and the
sub-Gaussian norm of $\xi_{i}$.

\medskip
\noindent\textbf{Step 4: Combining the three contributions.}
The test MSE satisfies
\begin{align}
  \MSE_{\mathrm{test}}
  &= \E_{\bx}\bigl[(f(\bx;\theta^{*}) - g^{*}(\bx))^{2}\bigr] \notag \\
  &= \E_{\bx}\bigl[(S(\bx) + N(\bx) + R(\bx))^{2}\bigr] \notag \\
  &\leq \E_{\bx}\bigl[\bigl(\abs{S(\bx)} + \abs{N(\bx)}
         + \abs{R(\bx)}\bigr)^{2}\bigr] \notag \\
  &\leq 3\,\E_{\bx}\bigl[S(\bx)^{2}\bigr]
         + 3\,\E_{\bx}\bigl[N(\bx)^{2}\bigr]
         + 3\,\delta_{\mathrm{lin}}^{2}, \label{eq:triangle}
\end{align}
where the last inequality uses $(a+b+c)^{2} \leq 3(a^{2}+b^{2}+c^{2})$.

By the definition of signal bias (Definition~\ref{def:signal-bias}),
$\E_{\bx}[S(\bx)^{2}] = B_{\mathrm{signal}}^{2}$.  By~\eqref{eq:EZ}
and~\eqref{eq:Z-tail},
$\E_{\bx}[N(\bx)^{2}] \leq \sigma^{2}\tr(\bM) + C_{1}\sigma^{2}\norm{\bM}_{F}\sqrt{\log(1/\delta)}$
with probability at least $1 - \delta$.  Substituting into~\eqref{eq:triangle}
and absorbing the factor of $3$ into the constant $C$ yields the
bound~\eqref{eq:main-bound}.
\end{proof}

\begin{remark}\label{rem:kernel-special}
In kernel regression with a fixed kernel $K$, the Jacobian is
$\bJ = \bX$ (the data matrix in feature space), and the test--train
cross-correlation $\boldsymbol{\Gamma}_{jl}$ reduces to the kernel eigenvalue
structure.  In this case, $\bM$ is determined entirely by the kernel eigenvalues,
and $\tr(\bM)$ reduces to the effective-rank-based criteria studied
in~\cite{BLLT2020, TB2023, MSCBR2022}.  The two-parameter framework thus
strictly generalizes the single-parameter taxonomy.
\end{remark}

\subsection{Benign overfitting dichotomy}

We now specialize to the setting where the singular values and cross-correlation
exhibit power-law decay, which is common in practice and connects directly to
the taxonomy of Mallinar et al.~\cite{MSCBR2022}.

\begin{assumption}[Power-law spectral decay]\label{ass:power-law}
There exist constants $\alpha > 0$, $\beta > 0$, $c_{1}, c_{2} > 0$ such that
for all $j \geq 1$:
\begin{enumerate}
  \item The singular values of $\bJ$ satisfy
    $c_{1}\,j^{-\alpha} \leq \sigma_{j} \leq c_{1}^{-1}\,j^{-\alpha}$.
  \item The diagonal entries of $\boldsymbol{\Gamma}$ satisfy
    $c_{2}\,j^{-\beta} \leq \rho_{j} \leq c_{2}^{-1}\,j^{-\beta}$,
    where $\rho_{j} = \sqrt{\boldsymbol{\Gamma}_{jj}}$.
\end{enumerate}
\end{assumption}

\begin{theorem}[Benign Overfitting Dichotomy]\label{thm:dichotomy}
Under Assumption~\ref{ass:power-law}, the noise-variance term
$\sigma^{2}\tr(\bM)$ in the bound of Theorem~\ref{thm:main} exhibits the
following trichotomy:
\begin{enumerate}
  \item \textbf{Benign regime} ($\beta > \alpha + \tfrac{1}{2}$):
    $\tr(\bM) < \infty$, so the noise contribution to the test MSE is bounded.
  \item \textbf{Catastrophic regime} ($\beta < \alpha + \tfrac{1}{2}$):
    $\tr(\bM) = \infty$; equivalently, the tail sum $\sum_{j > k}\rho_{j}^{2}/\sigma_{j}^{2}$
    diverges for every finite $k$, and the noise contribution is unbounded.
  \item \textbf{Tempered regime} ($\beta = \alpha + \tfrac{1}{2}$):
    $\tr(\bM)$ diverges logarithmically; specifically,
    $\sum_{j=1}^{J} \rho_{j}^{2}/\sigma_{j}^{2} = \Theta(\log J)$ as
    $J \to \infty$.
\end{enumerate}
\end{theorem}

\begin{proof}
We analyze the trace of $\bM$ by computing the diagonal entries.  Since
$\bM = \boldsymbol{\Sigma}^{-1}\boldsymbol{\Gamma}\boldsymbol{\Sigma}^{-1}$
and $\boldsymbol{\Sigma}$ is diagonal, the diagonal entries of $\bM$ are
\begin{equation}\label{eq:M-diag}
  \bM_{jj} = \frac{\boldsymbol{\Gamma}_{jj}}{\sigma_{j}^{2}}
            = \frac{\rho_{j}^{2}}{\sigma_{j}^{2}}.
\end{equation}
Therefore,
\begin{equation}\label{eq:trM}
  \tr(\bM) = \sum_{j=1}^{n} \frac{\rho_{j}^{2}}{\sigma_{j}^{2}}.
\end{equation}
Since $\boldsymbol{\Sigma}^{-1}$ is diagonal, we have
$\tr(\boldsymbol{\Sigma}^{-1}\boldsymbol{\Gamma}\boldsymbol{\Sigma}^{-1})
= \sum_{j} \sigma_{j}^{-2}\boldsymbol{\Gamma}_{jj}
= \sum_{j} \rho_{j}^{2}/\sigma_{j}^{2}$,
using the identity $\tr(\mathbf{D}\mathbf{A}\mathbf{D}) = \sum_{j} D_{jj}^{2} A_{jj}$ for diagonal $\mathbf{D}$.  Note that the off-diagonal entries of $\boldsymbol{\Gamma}$ do not contribute to $\tr(\bM)$.

Under Assumption~\ref{ass:power-law}, each summand satisfies
\begin{equation}\label{eq:summand}
  \frac{\rho_{j}^{2}}{\sigma_{j}^{2}}
  \asymp \frac{j^{-2\beta}}{j^{-2\alpha}}
  = j^{2(\alpha - \beta)},
\end{equation}
where $\asymp$ denotes equality up to positive multiplicative constants
depending on $c_{1}$ and $c_{2}$.  The tail sum starting from any fixed index
$k$ is therefore
\begin{equation}\label{eq:tail}
  L_{\mathrm{tail}} = \sum_{j > k} \frac{\rho_{j}^{2}}{\sigma_{j}^{2}}
  \asymp \sum_{j > k} j^{2(\alpha - \beta)}.
\end{equation}

We now apply the integral comparison test.  The series
$\sum_{j=1}^{\infty} j^{2(\alpha - \beta)}$ converges if and only if the
exponent satisfies $2(\alpha - \beta) < -1$, that is, $\beta > \alpha + \tfrac{1}{2}$.

\medskip
\noindent\textbf{Case 1:} $\beta > \alpha + \tfrac{1}{2}$.  Then
$2(\alpha - \beta) < -1$, so the series $\sum_{j} j^{2(\alpha-\beta)}$
converges.  By the comparison bounds from Assumption~\ref{ass:power-law},
$\tr(\bM) = \sum_{j} \rho_{j}^{2}/\sigma_{j}^{2}$ is bounded above by
$c_{2}^{-2}\,c_{1}^{2} \sum_{j} j^{2(\alpha - \beta)} < \infty$.

\medskip
\noindent\textbf{Case 2:} $\beta < \alpha + \tfrac{1}{2}$.  Then
$2(\alpha - \beta) > -1$, and the series diverges.  By the lower bound in
Assumption~\ref{ass:power-law},
$\tr(\bM) \geq c_{2}^{2}\,c_{1}^{-2} \sum_{j} j^{2(\alpha-\beta)} = \infty$.

\medskip
\noindent\textbf{Case 3:} $\beta = \alpha + \tfrac{1}{2}$.  Then
$2(\alpha - \beta) = -1$, and the series becomes the harmonic series:
\begin{equation}
  \sum_{j=1}^{J} j^{-1} = \Theta(\log J).
\end{equation}
Therefore $\sum_{j=1}^{J} \rho_{j}^{2}/\sigma_{j}^{2} = \Theta(\log J)$,
confirming logarithmic divergence.

This completes the proof of the trichotomy.
\end{proof}

\begin{remark}\label{rem:two-vs-one}
The key distinction from the single-parameter taxonomy of
Mallinar et al.~\cite{MSCBR2022} is the following.  In that framework, both the
noise absorption and noise expression are controlled by the same eigenvalue
sequence $\{\lambda_{j}\}$ of the kernel, leading to a condition that depends
only on the decay rate of $\{\lambda_{j}\}$.  Our condition
$\beta > \alpha + \tfrac{1}{2}$ involves \emph{two independent} rates: $\alpha$
controls how efficiently the interpolation absorbs noise (via $\sigma_{j}$), and
$\beta$ controls how visible the noise is at test points (via $\rho_{j}$).  The
condition $\beta > \alpha + \tfrac{1}{2}$ captures the requirement that noise
directions must be \emph{invisible at test points}, not merely that noise be
small in magnitude.  After feature learning, $\beta$ can increase (the learned
representation suppresses noise visibility at test points) while $\alpha$
remains fixed, potentially converting catastrophic overfitting into benign
overfitting---a transition that cannot be explained by any single-parameter
criterion.
\end{remark}

%% ============================================================
\section{Application to Teacher--Student Networks}\label{sec:app}
%% ============================================================

We demonstrate the utility of the $(\alpha, \beta)$ framework by analyzing a
concrete setting: two-layer ReLU teacher--student networks trained by gradient
flow.  The following result shows that gradient flow achieves spectral separation
in the Jacobian, which in turn guarantees benign generalization through our
Theorem~\ref{thm:main}.

\subsection{Setting}

Consider the following teacher--student configuration.

\begin{itemize}
  \item \textbf{Teacher network.}  The ground-truth function is a two-layer
    ReLU network with $k$ hidden neurons:
    \begin{equation}\label{eq:teacher}
      f^{*}(\bx) = \sum_{j=1}^{k} a_{j}^{*}\,\relu(\bw_{j}^{*\top}\bx),
    \end{equation}
    where $\bw_{1}^{*}, \ldots, \bw_{k}^{*} \in \R^{d}$ are orthonormal vectors
    and $a_{1}^{*}, \ldots, a_{k}^{*} \in \R$ are the second-layer weights.

  \item \textbf{Student network.}  The student is a two-layer ReLU network with
    $p \gg k$ hidden neurons:
    \begin{equation}\label{eq:student}
      f_{\theta}(\bx) = \frac{1}{\sqrt{p}}
        \sum_{j=1}^{p} a_{j}\,\relu(\bw_{j}^{\top}\bx),
    \end{equation}
    where $\theta = (a_{1}, \bw_{1}, \ldots, a_{p}, \bw_{p})$ collects all
    trainable parameters.

  \item \textbf{Training.}  The student is trained by population gradient flow
    (gradient flow on the expected loss) starting from an initialization scale
    $\alpha_{\mathrm{init}} \to 0$ (the ``rich'' or ``feature-learning'' regime).
    The training inputs $\bx_{i}$ are drawn i.i.d.\ from $\mathcal{N}(\mathbf{0}, \bI_{d})$.
\end{itemize}

\subsection{Spectral separation result}

\begin{proposition}[Jacobian Spectral Separation]\label{prop:teacher-student}
In the setting described above, suppose $n$ is sufficiently large relative to
$k$ and $d$, and $p \gg k(d+1)$.  Then at a global minimizer $\theta^{*}$ of the
population risk reached by gradient flow, the following spectral separation holds
for the training Jacobian $\bJ = \bJ(\theta^{*})$:
\begin{enumerate}
  \item For $i \leq k(d+1)$: $\sigma_{i}(\bJ) \geq c\sqrt{n}$, where $c > 0$
    is a constant depending on the teacher weights.
  \item For $i > k(d+1)$: $\sigma_{i}(\bJ) \leq C\alpha_{\mathrm{init}}$,
    where $C > 0$ is an absolute constant.
\end{enumerate}
Consequently, the effective rank satisfies $r_{\mathrm{eff}} \leq k(d+1)$, and
\begin{equation}\label{eq:gen-bound-ts}
  \MSE_{\mathrm{test}}
  \leq O\!\left(\sqrt{\frac{k\,d\,\log n}{n}}\right).
\end{equation}
\end{proposition}

\begin{proof}[Proof outline]
The argument combines results from Boursier and Flammarion~\cite{BF2022} with
our Theorem~\ref{thm:main}.  We proceed in four steps.

\medskip
\noindent\textbf{Step 1: Gradient flow alignment (cited result).}
By the main theorem of Boursier and Flammarion~\cite{BF2022}, in the rich regime
($\alpha_{\mathrm{init}} \to 0$), population gradient flow drives the student
network to a $k$-sparse aligned solution: at the global minimizer, exactly $k$
student neurons have their first-layer weights aligned with the teacher
directions $\bw_{1}^{*}, \ldots, \bw_{k}^{*}$, with second-layer weights
matching $a_{1}^{*}, \ldots, a_{k}^{*}$.  The remaining $p - k$ neurons have
weights of order $O(\alpha_{\mathrm{init}})$.

\medskip
\noindent\textbf{Step 2: Active block of the Jacobian.}
The $k$ aligned neurons contribute an ``active block'' $\bJ_{\mathrm{act}}$ to the
training Jacobian.  For each aligned neuron $j$ ($j = 1, \ldots, k$), the
gradient $\nabla_{\bw_{j}} f_{\theta^{*}}(\bx_{i})$ involves the indicator
$\mathbf{1}[\bw_{j}^{*\top}\bx_{i} > 0]$ and the input $\bx_{i}$.  Since
$\bw_{1}^{*}, \ldots, \bw_{k}^{*}$ are orthonormal and the inputs are Gaussian,
by a general-position argument (the probability that any collection of
$n$ Gaussian vectors lies in a degenerate configuration with respect to
$k$ hyperplanes is zero), the active block has rank $k(d+1)$: each neuron
contributes $d$ degrees of freedom from the first-layer weight gradient and $1$
from the second-layer weight, and the $k$ neurons contribute independently due
to orthonormality.  For $n$ sufficiently large relative to $k(d+1)$, the minimum
singular value of $\bJ_{\mathrm{act}}$ is at least $c\sqrt{n}$ by standard
random matrix concentration results for Gaussian designs.

\medskip
\noindent\textbf{Step 3: Dead block bound.}
The remaining $p - k$ neurons have weights of order
$O(\alpha_{\mathrm{init}})$ by Step~1.  Their contribution to the Jacobian,
denoted $\bJ_{\mathrm{dead}}$, satisfies
$\norm{\bJ_{\mathrm{dead}}}_{\mathrm{op}} = O(\alpha_{\mathrm{init}})$.
Each dead neuron has weights of order $O(\alpha_{\mathrm{init}})$, and the
operator norm of its rank-$(d+1)$ Jacobian block scales linearly with
$\alpha_{\mathrm{init}}$.  This means the dead neurons contribute negligibly to
the Jacobian spectrum.

\medskip
\noindent\textbf{Step 4: Spectral separation via Weyl's inequality.}
The full Jacobian is $\bJ = \bJ_{\mathrm{act}} + \bJ_{\mathrm{dead}}$.
By Weyl's inequality for singular values,
\begin{equation}
  \sigma_{i}(\bJ) \geq \sigma_{i}(\bJ_{\mathrm{act}})
    - \norm{\bJ_{\mathrm{dead}}}_{\mathrm{op}}.
\end{equation}
For $i \leq k(d+1)$: $\sigma_{i}(\bJ) \geq c\sqrt{n} - O(\alpha_{\mathrm{init}}) \geq c\sqrt{n}/2$ for $\alpha_{\mathrm{init}}$ sufficiently small.
For $i > k(d+1)$: $\sigma_{i}(\bJ_{\mathrm{act}}) = 0$ (since $\bJ_{\mathrm{act}}$ has rank at most $k(d+1)$), so $\sigma_{i}(\bJ) \leq \norm{\bJ_{\mathrm{dead}}}_{\mathrm{op}} = O(\alpha_{\mathrm{init}})$ by the reverse Weyl inequality.

This establishes the spectral gap.  The generalization bound~\eqref{eq:gen-bound-ts}
follows from Theorem~\ref{thm:main}: with effective rank
$r_{\mathrm{eff}} \leq k(d+1)$, the trace of $\bM$ is dominated by the
$k(d+1)$ large singular values, giving
$\tr(\bM) \leq k(d+1) \cdot O(1/n)$, and the standard $\sqrt{\log n / n}$
rate follows from the concentration term.
\end{proof}

\begin{remark}\label{rem:feature-learning}
This result illustrates the mechanism described in Remark~\ref{rem:two-vs-one}.
Before feature learning (i.e., at initialization with a random kernel), the
singular values of $\bJ$ would not exhibit spectral separation: all $n$ singular
values would be of comparable magnitude, and $\tr(\bM)$ would grow with $n$.
After gradient-flow training, the learned features concentrate the Jacobian
spectrum on $k(d+1)$ directions aligned with the teacher, effectively increasing
$\beta$ (the noise directions become invisible at test points) while $\alpha$
for the active directions remains unchanged.  This is the $(\alpha, \beta)$
decoupling in action.
\end{remark}

\begin{remark}\label{rem:proposition-status}
Proposition~\ref{prop:teacher-student} is labeled as a proposition rather than a
theorem because its proof relies on the gradient-flow alignment result
of~\cite{BF2022}, which is established under specific conditions (population
gradient flow, Gaussian inputs, orthonormal teacher weights, rich regime).
Extending this to finite-sample gradient descent with non-Gaussian inputs
remains an open problem.
\end{remark}

%% ============================================================
\section{Discussion}\label{sec:discussion}
%% ============================================================

\subsection{Connection to double descent}

The noise propagation operator $\bM$ provides a natural lens through which to
view the double-descent phenomenon~\cite{HMRT2022}.  Consider least-squares
regression with $n$ samples and $p$ features, and let $\gamma = p/n$ denote the
overparameterization ratio.

The trace $\tr(\bM)$ is dominated by the smallest singular values of the
Jacobian:
\begin{equation}\label{eq:trM-dom}
  \tr(\bM) = \sum_{j=1}^{n} \frac{\rho_{j}^{2}}{\sigma_{j}^{2}},
\end{equation}
and the sum is most sensitive to the indices $j$ where $\sigma_{j}$ is smallest.

In the underparameterized regime ($\gamma < 1$), the Marchenko--Pastur law for
random design matrices predicts that the smallest singular value of $\bJ$ is of
order $\sqrt{n}(1 - \sqrt{\gamma})$, and a calculation using the Marchenko--Pastur
density (see, e.g.,~\cite[Section~4]{HMRT2022}) yields
\begin{equation}
  \frac{\tr(\bM)}{n} \to \frac{(1 + \gamma)}{(1 - \gamma)^{3}}
  \qquad \text{as } n, p \to \infty \text{ with } p/n \to \gamma < 1.
\end{equation}
This expression diverges as $\gamma \to 1^{-}$, reflecting the peak of double
descent.

At the interpolation threshold ($\gamma = 1$), the smallest singular value
$\sigma_{\min}(\bJ) \to 0$ by the Marchenko--Pastur edge behavior, so
$\tr(\bM) \to \infty$.  This is the exact point where the interpolating solution
first exists, and it amplifies noise maximally.

In the overparameterized regime ($\gamma > 1$), the minimum-norm interpolant
selects the solution with the smallest parameter norm.  This implicit
regularization improves the effective conditioning of the problem: the relevant
singular values are those of the $n \times n$ matrix $\bJ\bJ^{\top}/p$, whose
smallest eigenvalue is of order $(\sqrt{\gamma} - 1)^{2}$, bounded away from
zero.  Consequently, $\tr(\bM)$ decreases as $\gamma$ increases beyond $1$,
completing the descent after the peak.

This analysis does not constitute a formal theorem, as it relies on asymptotic
random matrix theory results applied in a specific (linear, Gaussian design)
setting.  Extending it to the nonlinear setting covered by
Theorem~\ref{thm:main} requires controlling $\boldsymbol{\Gamma}$ in addition to
$\boldsymbol{\Sigma}$, which is an interesting direction for future work.

\subsection{Limitations}\label{subsec:limitations}

We identify the following limitations of our framework, which we state explicitly
to delineate the boundaries of our results.

\begin{enumerate}
  \item \textbf{Linearization residual is not quantified.}
    The bound in Theorem~\ref{thm:main} includes the term
    $\delta_{\mathrm{lin}}^{2}$, but we do not provide an explicit estimate of
    $\delta_{\mathrm{lin}}$ in terms of the network architecture or training
    procedure.  Our framework applies rigorously only when local linearization
    is a valid approximation (e.g., in the NTK regime or for wide networks).
    Quantifying $\delta_{\mathrm{lin}}$ for finite-width networks trained by
    gradient descent is a significant open problem.

  \item \textbf{Signal bias is not analyzed.}
    The term $B_{\mathrm{signal}}^{2}$ in~\eqref{eq:main-bound} is left as a
    given quantity.  Our analysis focuses entirely on the noise amplification
    mechanism; understanding $B_{\mathrm{signal}}^{2}$ requires a separate
    investigation of the approximation-theoretic properties of the learned
    features.

  \item \textbf{Population-level cross-correlation.}
    The matrix $\boldsymbol{\Gamma}$ is defined via a population expectation
    over test points (Definition~\ref{def:Gamma}).  In practice, one observes
    only finite test samples, and bounding the deviation between the empirical
    and population versions of $\boldsymbol{\Gamma}$ requires matrix
    concentration inequalities (e.g., matrix Bernstein bounds).  We defer this
    finite-sample analysis to future work.

  \item \textbf{Sub-Gaussian noise assumption.}
    The concentration step (Step~3 of the proof of Theorem~\ref{thm:main})
    relies on the Hanson--Wright inequality for sub-Gaussian random vectors.
    Extending the framework to heavy-tailed noise distributions would require
    alternative concentration tools, such as truncation arguments or the
    recently developed heavy-tailed Hanson--Wright inequalities.

  \item \textbf{Scope of the teacher--student result.}
    Proposition~\ref{prop:teacher-student} is established only for two-layer
    ReLU networks with orthonormal teacher weights, Gaussian inputs, and
    population gradient flow in the rich regime.  Extending the spectral
    separation result to deeper architectures, non-Gaussian inputs, or
    finite-sample stochastic gradient descent remains open.
\end{enumerate}

\subsection{Open problems}

We conclude by highlighting several directions for future investigation.

\begin{enumerate}
  \item \textbf{Quantifying $\delta_{\mathrm{lin}}$ for finite-width networks.}
    Can one bound the linearization residual for networks of moderate width
    trained by gradient descent, potentially using higher-order Taylor expansions
    or mean-field approximations?

  \item \textbf{Finite-sample $\boldsymbol{\Gamma}$ estimation.}
    What is the minimax rate for estimating $\tr(\bM)$ from a finite test sample,
    and can one construct practical estimators that predict the
    benign/tempered/catastrophic regime from data?

  \item \textbf{Feature-learning dynamics of $(\alpha, \beta)$.}
    Can one track the evolution of $\alpha$ and $\beta$ during gradient-descent
    training, and characterize the conditions under which training drives the
    system from the catastrophic to the benign regime?

  \item \textbf{Beyond power-law decay.}
    The dichotomy in Theorem~\ref{thm:dichotomy} assumes power-law decay.
    Extending the characterization to more general spectral profiles (e.g.,
    exponential decay, mixed regimes) would broaden the applicability of the
    framework.

  \item \textbf{Heavy-tailed and dependent noise.}
    Extending Theorem~\ref{thm:main} to accommodate heavy-tailed label noise
    or dependent noise structures (as arise in time-series or spatially
    correlated data) would significantly increase the practical relevance of the
    bounds.
\end{enumerate}

%% ============================================================
%% References
%% ============================================================

\subsection*{Acknowledgements}
The exploratory analysis and iterative refinement of the arguments in this
paper were conducted with the assistance of Claude (Anthropic).  All
mathematical statements and proofs have been verified by the author.

\begin{thebibliography}{10}

\bibitem{BHMM2019}
M.~Belkin, D.~Hsu, S.~Ma, and S.~Mandal.
\newblock Reconciling modern machine-learning practice and the classical
  bias--variance trade-off.
\newblock \emph{Proceedings of the National Academy of Sciences},
  116(32):15849--15854, 2019.

\bibitem{BLLT2020}
P.~L.~Bartlett, P.~M.~Long, G.~Lugosi, and A.~Tsigler.
\newblock Benign overfitting in linear regression.
\newblock \emph{Proceedings of the National Academy of Sciences},
  117(48):30063--30070, 2020.

\bibitem{BF2022}
E.~Boursier and N.~Flammarion.
\newblock Gradient flow dynamics of shallow {ReLU} networks for square loss and
  orthogonal inputs.
\newblock \emph{Advances in Neural Information Processing Systems}, 35, 2022.

\bibitem{HMRT2022}
T.~Hastie, A.~Montanari, S.~Rosset, and R.~J.~Tibshirani.
\newblock Surprises in high-dimensional ridgeless least squares interpolation.
\newblock \emph{Annals of Statistics}, 50(2):949--986, 2022.

\bibitem{JGH2018}
A.~Jacot, F.~Gabriel, and C.~Hongler.
\newblock Neural tangent kernel: convergence and generalization in neural
  networks.
\newblock \emph{Advances in Neural Information Processing Systems}, 31, 2018.

\bibitem{MM2022}
S.~Mei and A.~Montanari.
\newblock The generalization error of random features regression: Precise
  asymptotics and the double descent curve.
\newblock \emph{Communications on Pure and Applied Mathematics},
  75(4):667--766, 2022.

\bibitem{MSCBR2022}
N.~Mallinar, J.~B.~Simon, A.~Abedsoltan, P.~Pandit, M.~Belkin, and
  P.~Nakkiran.
\newblock Benign, tempered, or catastrophic: toward a refined taxonomy of
  overfitting.
\newblock \emph{Advances in Neural Information Processing Systems}, 35, 2022.

\bibitem{RZ2019}
A.~Rakhlin and X.~Zhai.
\newblock Consistency of interpolation with {L}aplace kernels is a
  high-dimensional phenomenon.
\newblock In \emph{Conference on Learning Theory (COLT)}, pages 2595--2623, 2019.

\bibitem{RV2013}
M.~Rudelson and R.~Vershynin.
\newblock Hanson--{W}right inequality and sub-{G}aussian concentration.
\newblock \emph{Electronic Communications in Probability}, 18:1--9, 2013.

\bibitem{TB2023}
A.~Tsigler and P.~L.~Bartlett.
\newblock Benign overfitting in ridge regression.
\newblock \emph{Journal of Machine Learning Research}, 24(123):1--76, 2023.

\end{thebibliography}

\end{document}
